% ============================================================
% CHAPITRE 24 — Design pour les réseaux sociaux
% ============================================================

\chapter{Design pour les réseaux sociaux}

\noindent
\textit{Dans le chapitre précédent, nous avons appris à créer des affiches et des flyers, supports de communication imprimés qui ont traversé les siècles. Mais aujourd'hui, une grande partie de la communication visuelle se fait sur les \emph{réseaux sociaux} : Instagram, Facebook, LinkedIn, TikTok, X (Twitter), Pinterest, YouTube. Ces plateformes sont devenues des canaux incontournables pour les marques, les entreprises, les associations, les créateurs. Elles imposent leurs propres codes, leurs propres formats, leurs propres contraintes.}

\vspace{1em}

\noindent
Le \terme{design pour les réseaux sociaux} est une discipline à part entière. Il ne suffit pas de transposer une affiche sur Instagram : il faut comprendre les spécificités de chaque plateforme, adapter les visuels aux formats, respecter les zones de sécurité, optimiser la lisibilité sur mobile, construire une série cohérente, organiser une campagne. Un visuel conçu pour un écran d'ordinateur ne fonctionne pas sur un smartphone ; une publication isolée n'a pas le même impact qu'une série cohérente ; un texte trop petit est illisible sur un fil d'actualité.

\vspace{1em}

\noindent
Ce chapitre a pour objectif de te faire maîtriser le design pour les réseaux sociaux. Nous commencerons par comprendre les objectifs de la communication numérique. Nous apprendrons ensuite à concevoir des publications statiques, à créer des visuels promotionnels, à construire une série de publications cohérentes, à adapter les compositions aux différents formats, à concevoir des couvertures et des bannières, à organiser une campagne visuelle, à préparer des exports adaptés aux plateformes, à éviter les problèmes de lisibilité sur téléphone, et à évaluer la cohérence d'une campagne. À la fin de ce chapitre, tu sauras créer des visuels efficaces et professionnels pour les réseaux sociaux.

% ============================================================
\section{Comprendre les objectifs de la communication numérique}
% ============================================================

\subsection{Qu'est-ce que la communication numérique ?}

\begin{quote}
\textbf{La communication numérique est l'ensemble des actions de communication menées sur les supports numériques : sites web, applications mobiles, réseaux sociaux, newsletters, publicités en ligne.}
\end{quote}

La communication numérique se distingue de la communication traditionnelle (imprimés, affichage) par :

\begin{itemize}
    \item \terme{La rapidité} : les messages sont diffusés instantanément.
    \item \terme{L'interactivité} : les utilisateurs peuvent réagir, commenter, partager.
    \item \terme{La mesure} : les résultats sont mesurables (statistiques, analytics).
    \item \terme{La personnalisation} : les contenus peuvent être adaptés aux utilisateurs.
    \item \terme{La mobilité} : les contenus sont consultés sur smartphone.
    \item \terme{La viralité} : les contenus peuvent se diffuser rapidement.
\end{itemize}

\subsection{Les objectifs de la communication numérique}

\subsubsection{Développer la notoriété}

L'objectif de \terme{notoriété} consiste à faire connaître une marque, un produit, un service, un événement. Il s'agit de toucher un large public, de gagner en visibilité, de se faire un nom.

\begin{itemize}
    \item \terme{Exemples} : campagne de lancement, publication virale, partenariat avec un influenceur.
    \item \terme{Indicateurs} : portée, impressions, mentions.
\end{itemize}

\subsubsection{Développer l'engagement}

L'objectif d'\terme{engagement} consiste à créer une relation avec le public : likes, commentaires, partages, sauvegardes. Il s'agit de susciter l'intérêt, de créer une communauté, de fidéliser.

\begin{itemize}
    \item \terme{Exemples} : contenu interactif, sondage, question, concours.
    \item \terme{Indicateurs} : taux d'engagement, nombre de commentaires, partages.
\end{itemize}

\subsubsection{Développer le trafic}

L'objectif de \terme{trafic} consiste à diriger les utilisateurs vers un site web, une boutique en ligne, une page d'inscription. Il s'agit de générer des visites, des clics, des conversions.

\begin{itemize}
    \item \terme{Exemples} : lien en bio, bouton d'appel à l'action, publicité.
    \item \terme{Indicateurs} : taux de clics, nombre de visites, taux de conversion.
\end{itemize}

\subsubsection{Développer les ventes}

L'objectif de \terme{ventes} consiste à inciter les utilisateurs à acheter un produit ou un service. Il s'agit de générer des commandes, des réservations, des abonnements.

\begin{itemize}
    \item \terme{Exemples} : promotion, offre limitée, code promo.
    \item \terme{Indicateurs} : nombre de ventes, chiffre d'affaires, panier moyen.
\end{itemize}

\subsubsection{Développer l'image}

L'objectif d'\terme{image} consiste à améliorer la perception d'une marque : valeurs, personnalité, positionnement. Il s'agit de construire une image cohérente, positive, mémorable.

\begin{itemize}
    \item \terme{Exemples} : communication corporate, behind the scenes, témoignages.
    \item \terme{Indicateurs} : sentiment, mentions, avis.
\end{itemize}

\subsubsection{Développer la fidélisation}

L'objectif de \terme{fidélisation} consiste à conserver les clients existants, à les faire revenir, à les transformer en ambassadeurs. Il s'agit de créer une relation durable, de récompenser la fidélité.

\begin{itemize}
    \item \terme{Exemples} : programme de fidélité, contenu exclusif, remerciement.
    \item \terme{Indicateurs} : taux de rétention, nombre de clients fidèles, recommandations.
\end{itemize}

\subsection{Les spécificités des réseaux sociaux}

\subsubsection{La temporalité}

Les réseaux sociaux sont \terme{éphémères} : un post est visible quelques heures, puis il est noyé dans le flux. Il faut donc publier régulièrement, saisir les moments forts, créer l'urgence.

\subsubsection{La concurrence}

Les réseaux sociaux sont \terme{saturés} : des millions de contenus sont publiés chaque jour. Il faut se distinguer, attirer l'attention, être mémorable.

\subsubsection{Le mobile}

Les réseaux sociaux sont consultés majoritairement sur \terme{smartphone}. Il faut donc concevoir des visuels adaptés aux petits écrans : lisibilité, contraste, format vertical.

\subsubsection{L'interaction}

Les réseaux sociaux sont \terme{interactifs} : les utilisateurs peuvent réagir, commenter, partager. Il faut donc créer du contenu engageant, susciter la conversation, répondre aux commentaires.

\subsubsection{La mesure}

Les réseaux sociaux sont \terme{mesurables} : les statistiques permettent de savoir ce qui fonctionne, ce qui ne fonctionne pas. Il faut donc analyser les résultats, ajuster la stratégie.

\subsection{Les bonnes pratiques}

\begin{itemize}
    \item \terme{Définir} un objectif clair.
    \item \terme{Choisir} les bonnes plateformes.
    \item \terme{Connaître} son public.
    \item \terme{Créer} du contenu adapté.
    \item \terme{Publier} régulièrement.
    \item \terme{Mesurer} les résultats.
    \item \terme{Ajuster} la stratégie.
\end{itemize}

\subsection*{Conclusion de la section}

La \terme{communication numérique} se distingue par sa rapidité, son interactivité, sa mesure, sa personnalisation, sa mobilité, sa viralité. Les objectifs sont : notoriété, engagement, trafic, ventes, image, fidélisation. Les réseaux sociaux ont des spécificités : temporalité, concurrence, mobile, interaction, mesure. Il est essentiel de définir un objectif clair, de connaître son public, de créer du contenu adapté, de mesurer les résultats.

Maintenant que nous connaissons les objectifs, nous allons, dans la section suivante, apprendre à \terme{concevoir des publications statiques}.

% ============================================================
\section{Concevoir des publications statiques}
% ============================================================

\subsection{Qu'est-ce qu'une publication statique ?}

\begin{quote}
\textbf{Une publication statique est un visuel fixe publié sur un réseau social : image, citation, infographie, photo, illustration. Elle s'oppose aux contenus animés (vidéos, GIF, stories animées).}
\end{quote}

Les publications statiques sont utilisées pour :

\begin{itemize}
    \item \terme{Présenter} un produit, un service.
    \item \terme{Partager} une citation, une information.
    \item \terme{Annoncer} un événement, une promotion.
    \item \terme{Montrer} les coulisses, l'équipe.
    \item \terme{Créer} de l'engagement.
\end{itemize}

\subsection{Les types de publications statiques}

\subsubsection{La publication visuelle}

La \terme{publication visuelle} met en avant une image forte : photographie, illustration, graphique. Elle est utilisée pour capter l'attention, créer une émotion.

\subsubsection{La publication informative}

La \terme{publication informative} transmet une information : conseil, astuce, actualité. Elle est utilisée pour informer, éduquer, engager.

\subsubsection{La publication promotionnelle}

La \terme{publication promotionnelle} promeut une offre : produit, service, promotion. Elle est utilisée pour vendre, convertir.

\subsubsection{La publication de citation}

La \terme{publication de citation} met en avant une phrase inspirante, un slogan, une citation. Elle est utilisée pour inspirer, engager, partager.

\subsubsection{La publication d'annonce}

La \terme{publication d'annonce} annonce un événement, un lancement, une nouveauté. Elle est utilisée pour informer, créer l'attente.

\subsection{Concevoir une publication statique}

\subsubsection{Étape 1 : Définir l'objectif}

\begin{itemize}
    \item \terme{Quel} est l'objectif de la publication ?
    \item \terme{Quel} est le public ?
    \item \terme{Quel} est le message ?
\end{itemize}

\subsubsection{Étape 2 : Choisir le format}

\begin{itemize}
    \item \terme{Carré} (1:1) : 1080 × 1080 px.
    \item \terme{Portrait} (4:5) : 1080 × 1350 px.
    \item \terme{Vertical} (9:16) : 1080 × 1920 px (story, reel).
    \item \terme{Paysage} (16:9) : 1200 × 675 px (LinkedIn, X).
\end{itemize}

\subsubsection{Étape 3 : Choisir le visuel}

\begin{itemize}
    \item \terme{Photographie} : pour un produit, une personne, une ambiance.
    \item \terme{Illustration} : pour un concept, une idée.
    \item \terme{Graphique} : pour une information, une donnée.
    \item \terme{Typographie} : pour une citation, un slogan.
\end{itemize}

\subsubsection{Étape 4 : Composer}

\begin{itemize}
    \item \terme{Placer} le visuel.
    \item \terme{Placer} le texte (court, lisible).
    \item \terme{Placer} le logo (discret).
    \item \terme{Placer} l'appel à l'action (si nécessaire).
\end{itemize}

\subsubsection{Étape 5 : Appliquer les couleurs}

\begin{itemize}
    \item \terme{Choisir} une palette cohérente.
    \item \terme{Vérifier} le contraste.
    \item \terme{Vérifier} la lisibilité.
\end{itemize}

\subsubsection{Étape 6 : Vérifier}

\begin{itemize}
    \item \terme{Vérifier} la lisibilité sur mobile.
    \item \terme{Vérifier} la cohérence.
    \item \terme{Vérifier} l'impact.
\end{itemize}

\subsection{Les bonnes pratiques}

\begin{itemize}
    \item \terme{Un seul message} par publication.
    \item \terme{Un visuel} fort.
    \item \terme{Un texte} court et lisible.
    \item \terme{Un appel à l'action} clair.
    \item \terme{Une palette} cohérente.
    \item \terme{Un logo} discret.
    \item \terme{Une lisibilité} sur mobile.
    \item \terme{Une cohérence} avec la charte.
\end{itemize}

\subsection*{Conclusion de la section}

Les \terme{publications statiques} sont des visuels fixes publiés sur les réseaux sociaux. Elles peuvent être visuelles, informatives, promotionnelles, de citation, ou d'annonce. La conception suit plusieurs étapes : définir l'objectif, choisir le format, choisir le visuel, composer, appliquer les couleurs, vérifier. Il est essentiel d'avoir un seul message, un visuel fort, un texte court, un appel à l'action clair, une lisibilité sur mobile.

Maintenant que nous savons concevoir des publications statiques, nous allons, dans la section suivante, apprendre à \terme{créer des visuels promotionnels}.

% ============================================================
\section{Créer des visuels promotionnels}
% ============================================================

\subsection{Qu'est-ce qu'un visuel promotionnel ?}

\begin{quote}
\textbf{Un visuel promotionnel est un visuel destiné à promouvoir un produit, un service, une offre, un événement. Il vise à attirer l'attention, à susciter l'intérêt, à déclencher l'action.}
\end{quote}

Les visuels promotionnels sont utilisés pour :

\begin{itemize}
    \item \terme{Annoncer} une promotion, une réduction.
    \item \terme{Lancer} un nouveau produit.
    \item \terme{Présenter} une offre limitée.
    \item \terme{Annoncer} un événement.
    \item \terme{Inciter} à l'achat.
\end{itemize}

\subsection{Les caractéristiques d'un visuel promotionnel efficace}

\subsubsection{Un impact visuel}

Un visuel promotionnel doit \terme{attirer l'attention} : couleurs vives, image forte, composition dynamique.

\subsubsection{Un message clair}

Un visuel promotionnel doit \terme{transmettre un message clair} : offre, avantage, bénéfice.

\subsubsection{Une urgence}

Un visuel promotionnel doit \terme{suggérer l'urgence} : date limite, stock limité, offre exceptionnelle.

\subsubsection{Un appel à l'action}

Un visuel promotionnel doit \terme{inciter à l'action} : bouton, flèche, indication.

\subsubsection{Une cohérence}

Un visuel promotionnel doit être \terme{cohérent} avec l'identité de la marque : logo, couleurs, typographies.

\subsection{Créer un visuel promotionnel}

\subsubsection{Étape 1 : Analyser l'offre}

\begin{itemize}
    \item \terme{Quel} est le produit ou service ?
    \item \terme{Quelle} est l'offre ?
    \item \terme{Quel} est le public ?
    \item \terme{Quel} est l'objectif ?
\end{itemize}

\subsubsection{Étape 2 : Choisir le visuel}

\begin{itemize}
    \item \terme{Choisir} une image du produit.
    \item \terme{Détourer} si nécessaire.
    \item \terme{Mettre} en valeur le produit.
\end{itemize}

\subsubsection{Étape 3 : Composer}

\begin{itemize}
    \item \terme{Placer} le produit.
    \item \terme{Placer} l'offre (prix, réduction).
    \item \terme{Placer} le message (accroche).
    \item \terme{Placer} l'appel à l'action.
    \item \terme{Placer} le logo.
\end{itemize}

\subsubsection{Étape 4 : Appliquer les couleurs}

\begin{itemize}
    \item \terme{Choisir} des couleurs vives, contrastées.
    \item \terme{Utiliser} la couleur pour attirer l'attention.
    \item \terme{Vérifier} la lisibilité.
\end{itemize}

\subsubsection{Étape 5 : Vérifier}

\begin{itemize}
    \item \terme{Vérifier} la lisibilité.
    \item \terme{Vérifier} l'impact.
    \item \terme{Vérifier} la cohérence.
\end{itemize}

\subsection{Les bonnes pratiques}

\begin{itemize}
    \item \terme{Un produit} visible.
    \item \terme{Une offre} claire.
    \item \terme{Un message} percutant.
    \item \terme{Un appel à l'action} visible.
    \item \terme{Une urgence} suggérée.
    \item \terme{Une palette} contrastée.
    \item \terme{Une cohérence} avec la marque.
    \item \terme{Une lisibilité} sur mobile.
\end{itemize}

\subsection*{Conclusion de la section}

Les \terme{visuels promotionnels} promeuvent un produit, un service, une offre. Ils doivent avoir un impact visuel, un message clair, une urgence, un appel à l'action, une cohérence. La création suit plusieurs étapes : analyser l'offre, choisir le visuel, composer, appliquer les couleurs, vérifier. Il est essentiel d'avoir un produit visible, une offre claire, un message percutant, un appel à l'action visible.

Maintenant que nous savons créer des visuels promotionnels, nous allons, dans la section suivante, apprendre à \terme{construire une série de publications cohérentes}.

% ============================================================
\section{Construire une série de publications cohérentes}
% ============================================================

\subsection{Pourquoi une série cohérente ?}

Une \terme{série de publications cohérentes} est essentielle pour :

\begin{itemize}
    \item \terme{Créer} une identité visuelle forte.
    \item \terme{Faciliter} la reconnaissance.
    \item \terme{Installer} une ambiance.
    \item \terme{Développer} une communauté.
    \item \terme{Marquer} les esprits.
    \item \terme{Optimiser} l'impact.
\end{itemize}

Une publication isolée a peu d'impact ; une série cohérente crée une dynamique, une attente, une habitude.

\subsection{Les principes de la cohérence}

\subsubsection{La cohérence visuelle}

Les publications doivent partager :

\begin{itemize}
    \item \terme{La} même palette de couleurs.
    \item \terme{Les} mêmes typographies.
    \item \terme{Le} même style de visuels.
    \item \terme{La} même composition (grille, marges).
    \item \terme{Le} même emplacement du logo.
\end{itemize}

\subsubsection{La cohérence éditoriale}

Les publications doivent partager :

\begin{itemize}
    \item \terme{Le} même ton (sérieux, ludique, inspirant).
    \item \terme{Le} même niveau de langage.
    \item \terme{Les} mêmes thématiques.
    \item \terme{Les} mêmes formats.
\end{itemize}

\subsubsection{La cohérence temporelle}

Les publications doivent être :

\begin{itemize}
    \item \terme{Régulières} : fréquence de publication.
    \item \terme{Planifiées} : calendrier éditorial.
    \item \terme{Coordonnées} : avec les autres canaux.
\end{itemize}

\subsection{Construire une série}

\subsubsection{Étape 1 : Définir le thème}

\begin{itemize}
    \item \terme{Quel} est le thème de la série ?
    \item \terme{Quel} est le message ?
    \item \terme{Quel} est le public ?
\end{itemize}

\subsubsection{Étape 2 : Définir la structure}

\begin{itemize}
    \item \terme{Définir} le nombre de publications.
    \item \terme{Définir} la fréquence.
    \item \terme{Définir} les formats (carré, portrait, story).
    \item \terme{Définir} les types (visuel, informatif, promotionnel).
\end{itemize}

\subsubsection{Étape 3 : Créer la charte de la série}

\begin{itemize}
    \item \terme{Définir} la palette de couleurs.
    \item \terme{Définir} les typographies.
    \item \terme{Définir} le style de visuels.
    \item \terme{Définir} la composition type.
    \item \terme{Définir} l'emplacement du logo.
\end{itemize}

\subsubsection{Étape 4 : Créer les publications}

\begin{itemize}
    \item \terme{Créer} chaque publication.
    \item \terme{Respecter} la charte.
    \item \terme{Vérifier} la cohérence.
\end{itemize}

\subsubsection{Étape 5 : Planifier la publication}

\begin{itemize}
    \item \terme{Créer} un calendrier éditorial.
    \item \terme{Planifier} les publications.
    \item \terme{Préparer} les textes (légendes, hashtags).
\end{itemize}

\subsubsection{Étape 6 : Publier et mesurer}

\begin{itemize}
    \item \terme{Publier} les visuels.
    \item \terme{Mesurer} les résultats.
    \item \terme{Ajuster} la stratégie.
\end{itemize}

\subsection{Les bonnes pratiques}

\begin{itemize}
    \item \terme{Définir} un thème clair.
    \item \terme{Définir} une structure.
    \item \terme{Créer} une charte de série.
    \item \terme{Respecter} la charte.
    \item \terme{Planifier} la publication.
    \item \terme{Vérifier} la cohérence.
    \item \terme{Mesurer} les résultats.
    \item \terme{Ajuster} la stratégie.
\end{itemize}

\subsection*{Conclusion de la section}

Une \terme{série de publications cohérentes} est essentielle pour créer une identité forte, faciliter la reconnaissance, installer une ambiance. Elle repose sur la cohérence visuelle, éditoriale, temporelle. La construction suit plusieurs étapes : définir le thème, la structure, la charte, créer les publications, planifier, publier, mesurer. Il est essentiel de définir un thème, une charte, de planifier, de vérifier la cohérence.

Maintenant que nous savons construire une série, nous allons, dans la section suivante, apprendre à \terme{adapter les compositions aux différents formats}.

% ============================================================
\section{Adapter les compositions aux différents formats}
% ============================================================

\subsection{Pourquoi adapter les compositions ?}

L'\terme{adaptation} des compositions aux différents formats est essentielle pour :

\begin{itemize}
    \item \terme{Communiquer} sur plusieurs plateformes.
    \item \terme{Toucher} plusieurs publics.
    \item \terme{Optimiser} l'impact.
    \item \terme{Respecter} les contraintes techniques.
    \item \terme{Assurer} la cohérence.
\end{itemize}

Chaque plateforme a ses formats, ses contraintes, ses codes. Une composition conçue pour un format ne fonctionne pas toujours dans un autre.

\subsection{Les formats par plateforme}

\subsubsection{Instagram}

\begin{itemize}
    \item \terme{Post carré} : 1080 × 1080 px (1:1).
    \item \terme{Post portrait} : 1080 × 1350 px (4:5).
    \item \terme{Story} : 1080 × 1920 px (9:16).
    \item \terme{Reel} : 1080 × 1920 px (9:16).
    \item \terme{Profil} : 320 × 320 px.
\end{itemize}

\subsubsection{Facebook}

\begin{itemize}
    \item \terme{Post} : 1200 × 630 px (1.91:1).
    \item \terme{Story} : 1080 × 1920 px (9:16).
    \item \terme{Couverture} : 851 × 315 px.
    \item \terme{Profil} : 170 × 170 px.
\end{itemize}

\subsubsection{LinkedIn}

\begin{itemize}
    \item \terme{Post} : 1200 × 627 px (1.91:1).
    \item \terme{Couverture} : 1584 × 396 px.
    \item \terme{Profil} : 400 × 400 px.
\end{itemize}

\subsubsection{X (Twitter)}

\begin{itemize}
    \item \terme{Post} : 1200 × 675 px (16:9).
    \item \terme{Couverture} : 1500 × 500 px.
    \item \terme{Profil} : 400 × 400 px.
\end{itemize}

\subsubsection{Pinterest}

\begin{itemize}
    \item \terme{Épingle} : 1000 × 1500 px (2:3).
    \item \terme{Profil} : 165 × 165 px.
\end{itemize}

\subsubsection{YouTube}

\begin{itemize}
    \item \terme{Miniature} : 1280 × 720 px (16:9).
    \item \terme{Bannière} : 2048 × 1152 px.
    \item \terme{Profil} : 800 × 800 px.
\end{itemize}

\subsection{Adapter une composition}

\subsubsection{Étape 1 : Créer la composition principale}

\begin{itemize}
    \item \terme{Créer} la composition dans un format de référence (par exemple, carré 1080 × 1080).
    \item \terme{Définir} les éléments essentiels.
    \item \terme{Définir} la hiérarchie.
\end{itemize}

\subsubsection{Étape 2 : Identifier les éléments à conserver}

\begin{itemize}
    \item \terme{Titre} (message principal).
    \item \terme{Visuel} (image, illustration).
    \item \terme{Logo}.
    \item \terme{Couleurs}.
    \item \terme{Typographies}.
\end{itemize}

\subsubsection{Étape 3 : Identifier les éléments à adapter}

\begin{itemize}
    \item \terme{Position} des éléments.
    \item \terme{Taille} des éléments.
    \item \terme{Marges}.
    \item \terme{Composition}.
\end{itemize}

\subsubsection{Étape 4 : Créer les déclinaisons}

\begin{itemize}
    \item \terme{Changer} le format.
    \item \terme{Repositionner} les éléments.
    \item \terme{Redimensionner} les éléments.
    \item \terme{Adapter} les marges.
\end{itemize}

\subsubsection{Étape 5 : Vérifier la cohérence}

\begin{itemize}
    \item \terme{Vérifier} que les déclinaisons sont cohérentes.
    \item \terme{Vérifier} la lisibilité.
    \item \terme{Vérifier} l'impact.
\end{itemize}

\subsubsection{Étape 6 : Exporter}

\begin{itemize}
    \item \terme{Exporter} dans le bon format.
    \item \terme{Exporter} dans la bonne résolution.
    \item \terme{Optimiser} le poids.
\end{itemize}

\subsection{Les bonnes pratiques}

\begin{itemize}
    \item \terme{Créer} une composition principale.
    \item \terme{Conserver} les éléments essentiels.
    \item \terme{Adapter} les éléments.
    \item \terme{Respecter} les proportions.
    \item \terme{Adapter} les marges.
    \item \terme{Vérifier} la cohérence.
    \item \terme{Vérifier} la lisibilité.
    \item \terme{Exporter} dans le bon format.
\end{itemize}

\subsection*{Conclusion de la section}

L'\terme{adaptation} des compositions aux différents formats est essentielle. Chaque plateforme a ses formats : Instagram (carré, portrait, story, reel), Facebook (post, story, couverture), LinkedIn (post, couverture), X (post, couverture), Pinterest (épingle), YouTube (miniature, bannière). La méthode consiste à créer la composition principale, à identifier les éléments à conserver et à adapter, à créer les déclinaisons, à vérifier la cohérence, à exporter. Il est essentiel de conserver les éléments essentiels, d'adapter les éléments, de vérifier la cohérence.

Maintenant que nous savons adapter les compositions, nous allons, dans la section suivante, apprendre à \terme{concevoir des couvertures et des bannières}.

% ============================================================
\section{Concevoir des couvertures et des bannières}
% ============================================================

\subsection{Qu'est-ce qu'une couverture ou une bannière ?}

\begin{quote}
\textbf{Une couverture (ou bannière) est un visuel de grand format qui orne le haut d'un profil, d'une page, d'une chaîne, ou qui sert de support publicitaire. Elle identifie, présente, promeut.}
\end{quote}

Les couvertures et bannières sont utilisées pour :

\begin{itemize}
    \item \terme{Identifier} un profil (Facebook, LinkedIn, YouTube, X).
    \item \terme{Présenter} une marque, un produit.
    \item \terme{Promouvoir} une offre, un événement.
    \item \terme{Créer} une ambiance.
    \item \terme{Différencier} le profil.
\end{itemize}

\subsection{Les types de couvertures et bannières}

\subsubsection{La couverture de profil}

La \terme{couverture de profil} orne le haut d'un profil (Facebook, LinkedIn, X, YouTube). Elle présente la marque, l'activité, les valeurs.

\subsubsection{La bannière publicitaire}

La \terme{bannière publicitaire} est un visuel promotionnel affiché sur des sites web, des applications, des réseaux sociaux. Elle promeut un produit, un service, une offre.

\subsubsection{La bannière de chaîne}

La \terme{bannière de chaîne} orne le haut d'une chaîne YouTube. Elle présente la chaîne, le contenu, la fréquence de publication.

\subsection{Les formats}

\begin{table}[h!]
\centering
\begin{tabular}{@{}p{3cm}p{4cm}p{6cm}@{}}
\toprule
\textbf{Plateforme} & \textbf{Type} & \textbf{Dimensions (px)} \\
\midrule
Facebook & Couverture & 851 × 315 \\
LinkedIn & Couverture & 1584 × 396 \\
X (Twitter) & Couverture & 1500 × 500 \\
YouTube & Bannière & 2048 × 1152 \\
Site web & Leaderboard & 728 × 90 \\
Site web & Medium Rectangle & 300 × 250 \\
Site web & Half Page & 300 × 600 \\
\bottomrule
\end{tabular}
\caption{Formats de couvertures et bannières.}
\end{table}

\subsection{Concevoir une couverture ou une bannière}

\subsubsection{Étape 1 : Analyser l'objectif}

\begin{itemize}
    \item \terme{Quel} est l'objectif ?
    \item \terme{Quel} est le public ?
    \item \terme{Quel} est le message ?
\end{itemize}

\subsubsection{Étape 2 : Choisir le visuel}

\begin{itemize}
    \item \terme{Choisir} une image forte, adaptée au format.
    \item \terme{Vérifier} la résolution.
    \item \terme{Vérifier} le cadrage.
\end{itemize}

\subsubsection{Étape 3 : Composer}

\begin{itemize}
    \item \terme{Placer} le visuel.
    \item \terme{Placer} le texte (court, lisible).
    \item \terme{Placer} le logo.
    \item \terme{Prévoir} une zone de sécurité (certains éléments peuvent être masqués par l'interface).
\end{itemize}

\subsubsection{Étape 4 : Appliquer les couleurs}

\begin{itemize}
    \item \terme{Choisir} une palette cohérente.
    \item \terme{Vérifier} le contraste.
    \item \terme{Vérifier} la lisibilité.
\end{itemize}

\subsubsection{Étape 5 : Vérifier}

\begin{itemize}
    \item \terme{Vérifier} la lisibilité.
    \item \terme{Vérifier} la cohérence.
    \item \terme{Vérifier} la zone de sécurité.
\end{itemize}

\subsection{Les zones de sécurité}

\subsubsection{Définition}

Les \terme{zones de sécurité} sont des zones où les éléments importants ne doivent pas être placés, car ils peuvent être masqués par l'interface (photo de profil, boutons, menus).

\subsubsection{Exemples}

\begin{itemize}
    \item \terme{Facebook} : la photo de profil masque le coin inférieur gauche.
    \item \terme{LinkedIn} : la photo de profil masque le coin inférieur gauche.
    \item \terme{X} : la photo de profil masque le coin inférieur gauche.
    \item \terme{YouTube} : la photo de profil et le nom de la chaîne masquent certaines zones.
\end{itemize}

\subsection{Les bonnes pratiques}

\begin{itemize}
    \item \terme{Analyser} l'objectif.
    \item \terme{Choisir} un visuel fort.
    \item \terme{Composer} avec soin.
    \item \terme{Prévoir} une zone de sécurité.
    \item \terme{Utiliser} une palette cohérente.
    \item \terme{Vérifier} la lisibilité.
    \item \terme{Vérifier} la cohérence.
    \item \terme{Tester} sur la plateforme.
\end{itemize}

\subsection*{Conclusion de la section}

Les \terme{couvertures et bannières} sont des visuels de grand format qui ornent le haut d'un profil, d'une page, d'une chaîne, ou qui servent de support publicitaire. Elles peuvent être des couvertures de profil, des bannières publicitaires, des bannières de chaîne. Les formats sont spécifiques à chaque plateforme. La conception suit plusieurs étapes : analyser l'objectif, choisir le visuel, composer, appliquer les couleurs, vérifier. Il est essentiel de prévoir une zone de sécurité, de vérifier la lisibilité, de tester sur la plateforme.

Maintenant que nous savons concevoir des couvertures et bannières, nous allons, dans la section suivante, apprendre à \terme{organiser une campagne visuelle}.

% ============================================================
\section{Organiser une campagne visuelle}
% ============================================================

\subsection{Qu'est-ce qu'une campagne visuelle ?}

\begin{quote}
\textbf{Une campagne visuelle est un ensemble coordonné de visuels diffusés sur plusieurs canaux (réseaux sociaux, site web, affichage, emailing), pendant une période donnée, autour d'un même message, d'un même objectif, d'une même identité.}
\end{quote}

Les campagnes visuelles sont utilisées pour :

\begin{itemize}
    \item \terme{Lancer} un produit.
    \item \terme{Promouvoir} une offre.
    \item \terme{Annoncer} un événement.
    \item \terme{Développer} la notoriété.
    \item \terme{Fidéliser} les clients.
\end{itemize}

\subsection{Les composantes d'une campagne}

\subsubsection{Le message}

Le \terme{message} est l'élément central : slogan, accroche, promesse. Il doit être clair, percutant, mémorable.

\subsubsection{Le visuel}

Le \terme{visuel} est l'élément fort : image, illustration, graphique. Il doit être impactant, cohérent, adapté à chaque support.

\subsubsection{Le calendrier}

Le \terme{calendrier} définit la durée, les étapes, les temps forts de la campagne.

\subsubsection{Les canaux}

Les \terme{canaux} sont les supports de diffusion : réseaux sociaux, site web, emailing, affichage, publicité en ligne.

\subsubsection{Le budget}

Le \terme{budget} définit les moyens alloués à la campagne.

\subsection{Organiser une campagne}

\subsubsection{Étape 1 : Définir l'objectif}

\begin{itemize}
    \item \terme{Quel} est l'objectif de la campagne ?
    \item \terme{Quel} est le public ?
    \item \terme{Quel} est le message ?
\end{itemize}

\subsubsection{Étape 2 : Définir le concept}

\begin{itemize}
    \item \terme{Définir} le concept créatif.
    \item \terme{Définir} le slogan.
    \item \terme{Définir} le visuel principal.
\end{itemize}

\subsubsection{Étape 3 : Décliner sur les canaux}

\begin{itemize}
    \item \terme{Adapter} le visuel à chaque canal.
    \item \terme{Respecter} les formats.
    \item \terme{Respecter} la cohérence.
\end{itemize}

\subsubsection{Étape 4 : Planifier}

\begin{itemize}
    \item \terme{Définir} le calendrier.
    \item \terme{Définir} les temps forts.
    \item \terme{Préparer} les contenus.
\end{itemize}

\subsubsection{Étape 5 : Lancer}

\begin{itemize}
    \item \terme{Publier} les visuels.
    \item \terme{Diffuser} les publicités.
    \item \terme{Envoyer} les emails.
\end{itemize}

\subsubsection{Étape 6 : Mesurer}

\begin{itemize}
    \item \terme{Mesurer} les résultats.
    \item \terme{Analyser} les performances.
    \item \terme{Ajuster} la stratégie.
\end{itemize}

\subsection{Les bonnes pratiques}

\begin{itemize}
    \item \terme{Définir} un objectif clair.
    \item \terme{Définir} un concept fort.
    \item \terme{Décliner} sur les canaux.
    \item \terme{Planifier} la campagne.
    \item \terme{Respecter} la cohérence.
    \item \terme{Mesurer} les résultats.
    \item \terme{Ajuster} la stratégie.
    \item \terme{Documenter} la campagne.
\end{itemize}

\subsection*{Conclusion de la section}

Une \terme{campagne visuelle} est un ensemble coordonné de visuels diffusés sur plusieurs canaux, autour d'un même message, d'un même objectif, d'une même identité. Elle repose sur un message, un visuel, un calendrier, des canaux, un budget. L'organisation suit plusieurs étapes : définir l'objectif, le concept, décliner, planifier, lancer, mesurer. Il est essentiel de définir un objectif, un concept, de décliner, de planifier, de mesurer.

Maintenant que nous savons organiser une campagne, nous allons, dans la section suivante, apprendre à \terme{préparer des exports adaptés aux plateformes}.

% ============================================================
\section{Préparer des exports adaptés aux plateformes}
% ============================================================

\subsection{Pourquoi préparer les exports ?}

La \terme{préparation des exports} est essentielle pour :

\begin{itemize}
    \item \terme{Respecter} les formats.
    \item \terme{Optimiser} le poids.
    \item \terme{Assurer} la qualité.
    \item \terme{Éviter} les problèmes d'affichage.
    \item \terme{Gagner} du temps.
\end{itemize}

Un export mal préparé peut entraîner une image floue, un poids excessif, un format inadapté.

\subsection{Les paramètres d'export}

\subsubsection{Le format}

\begin{itemize}
    \item \terme{JPEG} : pour les photos, les images sans transparence.
    \item \terme{PNG} : pour les images avec transparence.
    \item \terme{WebP} : format moderne, plus léger.
    \item \terme{MP4} : pour les vidéos.
    \item \terme{GIF} : pour les animations simples.
\end{itemize}

\subsubsection{La résolution}

\begin{itemize}
    \item \terme{72 ppi} : standard pour le web.
    \item \terme{144 ppi} : pour les écrans Retina.
\end{itemize}

\subsubsection{Les dimensions}

\begin{itemize}
    \item \terme{Respecter} les dimensions de la plateforme.
    \item \terme{Vérifier} le format (carré, portrait, paysage, story).
\end{itemize}

\subsubsection{Le profil colorimétrique}

\begin{itemize}
    \item \terme{sRGB} : standard pour le web.
    \item \terme{Éviter} le CMJN (pour l'impression).
\end{itemize}

\subsubsection{La compression}

\begin{itemize}
    \item \terme{JPEG} : qualité 60 à 80 (équilibre qualité/poids).
    \item \terme{PNG} : compression sans perte.
    \item \terme{WebP} : compression optimisée.
\end{itemize}

\subsubsection{Le poids}

\begin{itemize}
    \item \terme{Optimiser} le poids (moins de 200 Ko pour les bannières).
    \item \terme{Utiliser} des outils (TinyPNG, ImageOptim, Photoshop Export).
\end{itemize}

\subsection{Préparer les exports}

\subsubsection{Étape 1 : Vérifier les dimensions}

\begin{itemize}
    \item \terme{Vérifier} les dimensions.
    \item \terme{Vérifier} le format.
\end{itemize}

\subsubsection{Étape 2 : Choisir le format}

\begin{itemize}
    \item \terme{Choisir} le format adapté (JPEG, PNG, WebP).
    \item \terme{Choisir} la qualité.
\end{itemize}

\subsubsection{Étape 3 : Configurer la résolution}

\begin{itemize}
    \item \terme{Configurer} 72 ppi.
    \item \terme{Configurer} sRGB.
\end{itemize}

\subsubsection{Étape 4 : Exporter}

\begin{itemize}
    \item \terme{Exporter} dans Photoshop (Fichier > Exporter > Exporter sous).
    \item \terme{Exporter} dans Illustrator (Fichier > Exporter > Exporter pour les écrans).
    \item \terme{Utiliser} des outils (TinyPNG, ImageOptim).
\end{itemize}

\subsubsection{Étape 5 : Vérifier}

\begin{itemize}
    \item \terme{Vérifier} le poids.
    \item \terme{Vérifier} la qualité.
    \item \terme{Vérifier} le rendu.
\end{itemize}

\subsection{Les bonnes pratiques}

\begin{itemize}
    \item \terme{Respecter} les dimensions.
    \item \terme{Choisir} le bon format.
    \item \terme{Configurer} 72 ppi, sRGB.
    \item \terme{Optimiser} le poids.
    \item \terme{Vérifier} le rendu.
    \item \terme{Tester} sur la plateforme.
    \item \terme{Organiser} les fichiers.
    \item \terme{Documenter} les exports.
\end{itemize}

\subsection*{Conclusion de la section}

La \terme{préparation des exports} est essentielle pour respecter les formats, optimiser le poids, assurer la qualité. Les paramètres sont : format (JPEG, PNG, WebP), résolution (72 ppi), dimensions (adaptées à la plateforme), profil (sRGB), compression (qualité 60-80), poids (moins de 200 Ko). La préparation suit plusieurs étapes : vérifier les dimensions, choisir le format, configurer la résolution, exporter, vérifier. Il est essentiel de respecter les dimensions, de choisir le bon format, d'optimiser le poids.

Maintenant que nous savons préparer les exports, nous allons, dans la section suivante, apprendre à \terme{éviter les problèmes de lisibilité sur téléphone}.

% ============================================================
\section{Éviter les problèmes de lisibilité sur téléphone}
% ============================================================

\subsection{Pourquoi la lisibilité sur téléphone est-elle essentielle ?}

La \terme{lisibilité sur téléphone} est essentielle pour :

\begin{itemize}
    \item \terme{Toucher} la majorité du public (80\% des utilisateurs consultent sur mobile).
    \item \terme{Assurer} la compréhension du message.
    \item \terme{Faciliter} l'engagement.
    \item \terme{Éviter} les abandons.
    \item \terme{Optimiser} l'impact.
\end{itemize}

Un visuel illisible sur téléphone est un visuel raté, même s'il est beau sur ordinateur.

\subsection{Les problèmes de lisibilité}

\subsubsection{Le texte trop petit}

Sur un téléphone, le texte est vu en petit. Un texte trop petit est illisible.

\subsubsection{Le contraste insuffisant}

Un contraste insuffisant entre texte et fond rend le texte illisible, surtout en plein soleil.

\subsubsection{La surcharge}

Trop d'éléments, trop de texte, trop de couleurs : le visuel est confus, illisible.

\subsubsection{Le format inadapté}

Un format horizontal sur un écran vertical, ou un format carré sur une story, réduit la lisibilité.

\subsubsection{Les zones de sécurité ignorées}

Des éléments importants sont masqués par l'interface (photo de profil, boutons).

\subsection{Éviter les problèmes}

\subsubsection{Adapter la taille du texte}

\begin{itemize}
    \item \terme{Titre} : au moins 40 px.
    \item \terme{Sous-titre} : au moins 24 px.
    \item \terme{Texte courant} : au moins 16 px.
    \item \terme{Mentions} : au moins 12 px.
\end{itemize}

\subsubsection{Assurer le contraste}

\begin{itemize}
    \item \terme{Contraste} entre texte et fond (au moins 4.5:1).
    \item \terme{Contraste} entre éléments.
    \item \terme{Éviter} les textes clairs sur fond clair.
\end{itemize}

\subsubsection{Simplifier}

\begin{itemize}
    \item \terme{Un} seul message.
    \item \terme{Peu} de texte.
    \item \terme{Peu} d'éléments.
    \item \terme{Une} composition claire.
\end{itemize}

\subsubsection{Adapter le format}

\begin{itemize}
    \item \terme{Format} vertical pour les stories.
    \item \terme{Format} carré ou portrait pour les posts.
    \item \terme{Format} horizontal pour les bannières.
\end{itemize}

\subsubsection{Respecter les zones de sécurité}

\begin{itemize}
    \item \terme{Prévoir} une zone de sécurité.
    \item \terme{Vérifier} sur la plateforme.
    \item \terme{Éviter} de placer des éléments importants dans les zones masquées.
\end{itemize}

\subsection{Tester la lisibilité}

\subsubsection{Test sur téléphone}

\begin{itemize}
    \item \terme{Ouvrir} le visuel sur un téléphone.
    \item \terme{Vérifier} la lisibilité.
    \item \terme{Vérifier} le contraste.
    \item \terme{Vérifier} la taille du texte.
\end{itemize}

\subsubsection{Test sur différents appareils}

\begin{itemize}
    \item \terme{Tester} sur iPhone.
    \item \terme{Tester} sur Android.
    \item \terme{Tester} sur tablette.
\end{itemize}

\subsubsection{Test avec des personnes}

\begin{itemize}
    \item \terme{Montrer} le visuel à d'autres personnes.
    \item \terme{Demander} leur avis.
    \item \terme{Vérifier} la compréhension.
\end{itemize}

\subsection{Les bonnes pratiques}

\begin{itemize}
    \item \terme{Adapter} la taille du texte.
    \item \terme{Assurer} le contraste.
    \item \terme{Simplifier} la composition.
    \item \terme{Adapter} le format.
    \item \terme{Respecter} les zones de sécurité.
    \item \terme{Tester} sur téléphone.
    \item \terme{Tester} sur différents appareils.
    \item \terme{Demander} un avis.
\end{itemize}

\subsection*{Conclusion de la section}

La \terme{lisibilité sur téléphone} est essentielle pour toucher la majorité du public. Les problèmes sont : texte trop petit, contraste insuffisant, surcharge, format inadapté, zones de sécurité ignorées. Pour les éviter : adapter la taille du texte (au moins 40 px pour le titre), assurer le contraste, simplifier, adapter le format, respecter les zones de sécurité. Il est essentiel de tester sur téléphone, sur différents appareils, de demander un avis.

Maintenant que nous savons éviter les problèmes de lisibilité, nous allons, dans la dernière section de ce chapitre, apprendre à \terme{évaluer la cohérence d'une campagne}.

% ============================================================
\section{Évaluer la cohérence d'une campagne}
% ============================================================

\subsection{Pourquoi évaluer la cohérence ?}

L'\terme{évaluation de la cohérence} est essentielle pour :

\begin{itemize}
    \item \terme{Vérifier} que la campagne est cohérente.
    \item \terme{Assurer} la reconnaissance.
    \item \terme{Optimiser} l'impact.
    \item \terme{Corriger} les incohérences.
    \item \terme{Améliorer} les futures campagnes.
\end{itemize}

Une campagne incohérente brouille le message, nuit à la marque, réduit l'efficacité.

\subsection{Les critères de cohérence}

\subsubsection{La cohérence visuelle}

\begin{itemize}
    \item \terme{Même} palette de couleurs.
    \item \terme{Mêmes} typographies.
    \item \terme{Même} style de visuels.
    \item \terme{Même} composition.
    \item \terme{Même} emplacement du logo.
\end{itemize}

\subsubsection{La cohérence éditoriale}

\begin{itemize}
    \item \terme{Même} ton.
    \item \terme{Même} niveau de langage.
    \item \terme{Mêmes} messages.
    \item \terme{Mêmes} hashtags.
\end{itemize}

\subsubsection{La cohérence temporelle}

\begin{itemize}
    \item \terme{Même} calendrier.
    \item \terme{Même} fréquence.
    \item \terme{Mêmes} temps forts.
\end{itemize}

\subsubsection{La cohérence entre canaux}

\begin{itemize}
    \item \terme{Même} message.
    \item \terme{Même} visuel.
    \item \terme{Même} identité.
    \item \terme{Adaptation} aux formats.
\end{itemize}

\subsection{Évaluer la cohérence}

\subsubsection{Étape 1 : Rassembler les visuels}

\begin{itemize}
    \item \terme{Rassembler} tous les visuels de la campagne.
    \item \terme{Les} organiser par canal.
    \item \terme{Les} organiser par format.
\end{itemize}

\subsubsection{Étape 2 : Analyser la cohérence}

\begin{itemize}
    \item \terme{Comparer} les visuels.
    \item \terme{Vérifier} la palette.
    \item \terme{Vérifier} les typographies.
    \item \terme{Vérifier} les messages.
    \item \terme{Vérifier} le ton.
\end{itemize}

\subsubsection{Étape 3 : Identifier les incohérences}

\begin{itemize}
    \item \terme{Identifier} les écarts.
    \item \terme{Identifier} les incohérences.
    \item \terme{Identifier} les éléments à corriger.
\end{itemize}

\subsubsection{Étape 4 : Corriger}

\begin{itemize}
    \item \terme{Corriger} les incohérences.
    \item \terme{Aligner} les visuels.
    \item \terme{Vérifier} la cohérence.
\end{itemize}

\subsubsection{Étape 5 : Valider}

\begin{itemize}
    \item \terme{Valider} la cohérence.
    \item \terme{Documenter} les règles.
    \item \terme{Préparer} les futures campagnes.
\end{itemize}

\subsection{Les outils d'évaluation}

\subsubsection{Le moodboard}

Le \terme{moodboard} permet de rassembler les références visuelles et de vérifier la cohérence.

\subsubsection{La charte de campagne}

La \terme{charte de campagne} définit les règles visuelles et éditoriales de la campagne.

\subsubsection{Le calendrier éditorial}

Le \terme{calendrier éditorial} permet de vérifier la cohérence temporelle.

\subsubsection{Les indicateurs}

Les \terme{indicateurs} (taux d'engagement, portée, conversions) permettent de mesurer l'efficacité.

\subsection{Les bonnes pratiques}

\begin{itemize}
    \item \terme{Rassembler} les visuels.
    \item \terme{Analyser} la cohérence.
    \item \terme{Identifier} les incohérences.
    \item \terme{Corriger} les incohérences.
    \item \terme{Valider} la cohérence.
    \item \terme{Documenter} les règles.
    \item \terme{Préparer} les futures campagnes.
    \item \terme{Mesurer} l'efficacité.
\end{itemize}

\subsection*{Conclusion de la section}

L'\terme{évaluation de la cohérence} est essentielle pour vérifier que la campagne est cohérente, assurer la reconnaissance, optimiser l'impact. Les critères sont : cohérence visuelle, éditoriale, temporelle, entre canaux. L'évaluation suit plusieurs étapes : rassembler les visuels, analyser la cohérence, identifier les incohérences, corriger, valider. Il est essentiel de rassembler, d'analyser, de corriger, de valider, de documenter.

% ============================================================
\section*{Conclusion du chapitre}
\addcontentsline{toc}{section}{Conclusion du chapitre}
% ============================================================

Nous voici au terme du chapitre 24. Nous avons appris à maîtriser le design pour les réseaux sociaux :

\begin{itemize}
    \item \terme{Objectifs de la communication numérique} : notoriété, engagement, trafic, ventes, image, fidélisation.
    \item \terme{Publications statiques} : types (visuelle, informative, promotionnelle, citation, annonce), conception.
    \item \terme{Visuels promotionnels} : impact, message, urgence, appel à l'action, cohérence.
    \item \terme{Séries cohérentes} : cohérence visuelle, éditoriale, temporelle ; construction.
    \item \terme{Adaptation aux formats} : formats par plateforme, méthode.
    \item \terme{Couvertures et bannières} : types, formats, zones de sécurité, conception.
    \item \terme{Campagnes visuelles} : message, visuel, calendrier, canaux, budget ; organisation.
    \item \terme{Exports adaptés} : format, résolution, dimensions, profil, compression, poids.
    \item \terme{Lisibilité sur téléphone} : problèmes, solutions, tests.
    \item \terme{Cohérence de campagne} : critères, évaluation, correction.
\end{itemize}

Le design pour les réseaux sociaux est une discipline à part entière. Il ne suffit pas de transposer une affiche sur Instagram : il faut comprendre les spécificités de chaque plateforme, adapter les visuels aux formats, respecter les zones de sécurité, optimiser la lisibilité sur mobile, construire une série cohérente, organiser une campagne. Un visuel conçu pour un écran d'ordinateur ne fonctionne pas sur un smartphone ; une publication isolée n'a pas le même impact qu'une série cohérente ; un texte trop petit est illisible sur un fil d'actualité.

Ce chapitre t'a donné les bases du design pour les réseaux sociaux. Mais la maîtrise vient avec la pratique. N'hésite pas à t'entraîner sur des projets variés : posts, stories, couvertures, campagnes. Chaque projet est une occasion de progresser.

Dans le chapitre suivant, nous aborderons le \terme{branding et les systèmes d'identité visuelle}. Tu apprendras à comprendre la différence entre logo, identité visuelle et branding, à effectuer des recherches et un positionnement visuel, à créer un logo et ses variantes, à choisir les couleurs de marque, à définir le système typographique, à créer des éléments graphiques secondaires, à appliquer l'identité sur les supports, à élaborer une charte graphique, et à présenter un projet d'identité visuelle. Ces compétences te permettront de créer des marques fortes et cohérentes.

\vspace{2em}

\begin{center}
\textit{Fin du chapitre 24 --- Design pour les réseaux sociaux.}

\vspace{1em}

\textbf{Prochain chapitre : Chapitre 25 --- Branding et systèmes d'identité visuelle.}
\end{center}